\motto{The mathematician's patterns, like the painter's or the poet's must be beautiful; the ideas, like the colors or the words must fit together in a harmonious way. Beauty is the first test: There is no permanent place in this world for ugly mathematics.

--- Godfrey Harold Hardy\footnote{Godfrey Harold Hardy (1877--1947) was a British mathematician famous for his work in number theory and mathematical analysis. Apart from his research, Hardy was a strong advocate for pure mathematics and believed that mathematics should be pursued for its own beauty, not just for practical use.

He remained lifelong unmarried and dedicated much of his life entirely to mathematics, fitting into the common stereotype of a mathematician.}}
\chapter{The theory of b-divisors}
\label{chap:bdiv}

The purpose of this chapter is to translate the singularity theory developed in the preceding chapters into the language of b-divisors. A b-divisor records, on all proper bimeromorphic models of $X$, the cohomological data attached to a singular object.

Given a closed positive $(1,1)$-current $T$ on $X$, we associate to it a b-divisor $\mathbb D(T)$ by taking, on each modification $\pi\colon Y\to X$, the cohomology class of the regular part of $\pi^*T$ with respect to Siu's decomposition. This construction removes divisorial components and keeps track of the moving positivity of $T$ on every model. The resulting b-divisor is nef in a suitable sense, and its volume agrees with the volume of $T$.

The first part of the chapter develops this construction and studies its basic properties. We prove that $d_S$-convergence of currents implies convergence of the associated b-divisors, and we compare the $P$- and $\mathcal I$-orders through their b-divisorial data. This gives a birational way to read some of the singularity information encoded by envelopes and multiplier ideals.

The second part relates b-divisors to $\mathcal I$-good currents. In the big case, $\mathcal I$-good non-divisorial currents give rise to big and nef b-divisors, and conversely big and nef b-divisors can be represented by such currents.

The final part defines mixed intersection numbers of nef b-divisors. For big and nef b-divisors this is done through the corresponding $\mathcal I$-good currents and the mixed volume formalism of \cref{chap:Igood}; the general nef case is obtained by K\"ahler perturbation. We prove symmetry, multilinearity, monotonicity, continuity under decreasing nets, and upper semicontinuity. These results provide the b-divisorial counterpart of the mixed volume theory.


\section{The notion of b-divisors}\label{sec:bdivnotion}

The b-divisors defined in this section are sometimes called b-divisor classes. We always omit the word \emph{classes} to save space.

Let $X$ be a connected compact Kähler manifold of dimension $n$.

Let us recall the following elementary result regarding how cohomology behaves under blow-up.
\begin{proposition}\label{prop:blowup_cone}
    Let $\pi\colon Y\rightarrow X$ be a blow-up with \emph{connected} smooth center of codimension at least $2$ with exceptional divisor $E$. Then there is a natural identification
    \begin{equation}\label{eq:H11blowup}
    \mathrm{H}^{1,1}(Y,\mathbb{R})=\mathrm{H}^{1,1}(X,\mathbb{R})\oplus \mathbb{R}\{E\}.
    \end{equation}
\end{proposition}
See \cite{RYY19} for a much more general result.


\subsection{Nef b-divisors}
\label{subsec:nefbdiv}
\begin{definition}\label{def:bdiv:L38}
    A \emph{(Weil) b-divisor}\index{b-divisor} $\mathbb{D}$ over $X$ is an assignment $(\mathbb{D}_{\pi})_{\pi\colon Y\rightarrow X}$\index{$\mathbb{D}$}, where $\pi\colon Y\rightarrow X$ runs over all modifications of $X$ such that
    \begin{enumerate}
        \item\label{itm:chapter-bdiv:L43:1} $\mathbb{D}_{\pi}\in \mathrm{H}^{1,1}(Y,\mathbb{R})$;
        \item\label{itm:chapter-bdiv:L43:2} The classes are compatible under push-forwards: If $\pi'\colon Z\rightarrow X$ and $\pi\colon Y\rightarrow X$ are both modifications of $X$ and $\pi'$ dominates $\pi$ through $g\colon Z\rightarrow Y$, then $g_*\mathbb{D}_{\pi'}=\mathbb{D}_{\pi}$.
    \end{enumerate}
    We also write $\mathbb{D}_Y=\mathbb{D}_{\pi}$ if there is no risk of confusion.

    Given two Weil b-divisors $\mathbb{D}$ and $\mathbb{D}'$ over $X$, we say $\mathbb{D}\leq \mathbb{D}'$ if for each modification $\pi$ of $X$, we have $\mathbb{D}_{\pi}\leq \mathbb{D}'_{\pi}$. Recall that by definition, this means the class $\mathbb{D}'_{\pi}-\mathbb{D}_{\pi}$ is pseudo-effective.

    A Weil b-divisor with $\mathbb{D}\geq 0$ is called \emph{pseudo-effective}\index{b-divisor!pseudo-effective b-divisor}.
\end{definition}
The class $\mathbb{D}_X$ is called the \emph{root}\index{root!root of a b-divisor} of $\mathbb{D}$.
The set of Weil b-divisors over $X$ has the natural structure of a real vector space.


\begin{definition}\label{def:volb}
    The \emph{volume}\index{volume!volume of a b-divisor} of a pseudo-effective Weil b-divisor $\mathbb{D}$ over $X$ is
    \[
        \vol \mathbb{D}\coloneqq \lim_{\pi\colon Y\rightarrow X} \vol \mathbb{D}_Y.
    \]
    The right-hand side is a decreasing net due to \cref{prop:cones_modif2}, hence the limit always exists. The volume of a non-pseudo-effective Weil b-divisor is defined to be $0$.

    We say $\mathbb{D}$ is \emph{big}\index{b-divisor!big b-divisor} if $\vol \mathbb{D}>0$.
\end{definition}

\begin{lemma}\label{lma:volbdivnet}
    Let $(\mathbb{D}_i)_{i\in I}$ be a net of pseudo-effective b-divisors converging to a Weil b-divisor $\mathbb{D}$. Then
    \begin{equation}\label{eq:vollimsup}
    \varlimsup_{i\in I}\vol \mathbb{D}_i\leq \vol \mathbb{D}.
    \end{equation}
    If the net is decreasing, then
    \[
    \lim_{i\in I}\vol \mathbb{D}_i=\vol \mathbb{D}.
    \]
\end{lemma}
Here we say $(\mathbb{D}_i)_{i\in I}$ converges to $\mathbb{D}$ if for any modification $\pi\colon Y\rightarrow X$, we have $\mathbb{D}_{i,Y}\to \mathbb{D}_Y$ with respect to the Euclidean topology.
\begin{proof}
    First observe that $\mathbb{D}$ is pseudo-effective, since $\mathbb{D}_i\geq 0$ for all $i$.

    Let $\pi\colon Y\rightarrow X$ be a modification. Then
    \[
    \vol \mathbb{D}_Y=\lim_{i\in I} \vol \mathbb{D}_{i,Y}\geq \varlimsup_{i\in I} \vol \mathbb{D}_i.
    \]
    The inequality \eqref{eq:vollimsup} follows.
    As for the decreasing case, it suffices to observe that both sides of \eqref{eq:vollimsup} can be written as
    \[
    \inf_i \inf_{\pi\colon Y\rightarrow X} \vol \mathbb{D}_{i,Y}.
    \]
    This completes the proof.
\end{proof}

\begin{example}\label{ex:volbdivnetstrict}
    In general, equality need not hold in \eqref{eq:vollimsup}. Let $X=\mathbb{P}^2$, and let $H=c_1(\mathcal{O}_{\mathbb{P}^2}(1))$. Choose a sequence of points
    $p_i\in X$ converging to the generic point in the following sense: for every algebraic curve $C\subseteq X$, the set of indices $i$ with $p_i\in C$ is finite.

    Let
    \[
        \pi_i\colon X_i\rightarrow X
    \]
    be the blow-up of $X$ at $p_i$, with exceptional divisor $E_i$. Let $\mathbb{D}_i$ be the Cartier b-divisor realized on $X_i$ by the nef class
    \[
        \pi_i^*H-\{E_i\},
    \]
    and let $\mathbb{D}$ be the Cartier b-divisor realized on $X$ by $H$. Then $\mathbb{D}_i\to \mathbb{D}$.
    Indeed, for any fixed modification $\rho\colon Y\rightarrow X$, the morphism $\rho$ is an isomorphism over $p_i$ for all large enough $i$, hence the trace of $\mathbb{D}_i$ on $Y$ is $\rho^*H$ for all large enough $i$.

    On the other hand,
    \[
        \vol \mathbb{D}_i=\left(\pi_i^*H-\{E_i\}\right)^2=0
    \]
    for all $i$, whereas
    \[
        \vol \mathbb{D}=H^2=1.
    \]
    Thus the inequality in \eqref{eq:vollimsup} is strict for this sequence.
\end{example}

\begin{definition}\label{def:bdiv:L116}
    A \emph{Cartier b-divisor}\index{b-divisor!Cartier b-divisor} $\mathbb{D}$ over $X$ is a Weil b-divisor $\mathbb{D}$ over $X$ such that there exists a modification $\pi\colon Y\rightarrow X$ and a class $\alpha_Y\in \mathrm{H}^{1,1}(Y,\mathbb{R})$ so that for each $\pi'\colon Z\rightarrow X$ dominating $\pi$, the class $\mathbb{D}_Z$ is the pull-back of $\alpha_Y$. Any such $(\pi,\alpha_Y)$ is called a \emph{realization}\index{realization!realization of a Cartier b-divisor} of $\mathbb{D}$.
\end{definition}
By abuse of language, we also say $(Y,\alpha_Y)$ is a realization of $\mathbb{D}$.
The realization is not unique in general.


\begin{definition}\label{def:nefCar}
    A Cartier b-divisor $\mathbb{D}$ over $X$ is \emph{nef}\index{b-divisor!Cartier b-divisor!nef Cartier b-divisor} if there exists a realization $(\pi\colon Y\rightarrow X, \alpha_Y)$ of $\mathbb{D}$ such that $\alpha_Y$ is nef.
\end{definition}

\begin{definition}\label{def:nefb}
    A Weil b-divisor $\mathbb{D}$ over $X$ is \emph{nef}\index{b-divisor!nef b-divisor} if it is the limit of a net of nef Cartier b-divisors $(\mathbb{D}_i)_i$ over $X$.
\end{definition}
In other words, for each modification $\pi\colon Y\rightarrow X$, we have $\mathbb{D}_{i,Y}\to \mathbb{D}_Y$.

Note that thanks to \cref{prop:cones_modif}, each $\mathbb{D}_Y$ is necessarily modified nef, but it is not nef in general. The notion of modified nef classes is defined in \cref{def:mnclass}.

\emph{A priori}, for a Cartier b-divisor, nefness could mean either the notion defined by \cref{def:nefCar} or the one defined by \cref{def:nefb}. We will show in \cref{cor:nefbCequiv} that they are actually equivalent. Before that, by a nef Cartier b-divisor, we always mean nefness in the sense of \cref{def:nefCar}.

Our definition \cref{def:nefb} amounts to defining the set of Weil b-divisors as the closure of the set of Cartier b-divisors in $\varprojlim_{\pi}\mathrm{H}^{1,1}(Y,\mathbb{R})$ with respect to the projective limit topology. In particular, the limit of a converging net of nef b-divisors is still nef.

\subsection{The b-divisors of currents}\label{subsec:b-div_curr}

\begin{definition}\label{def:DT}
Let $T$ be a closed positive $(1,1)$-current on $X$. Given any modification $\pi\colon Y\rightarrow X$, we define
\begin{equation}
\mathbb{D}(T)_Y\coloneqq \left\{ \Reg \pi^*T \right\}\in \mathrm{H}^{1,1}(Y,\mathbb{R}).
\end{equation}
We call $\mathbb{D}(T)$\index{$\mathbb{D}(T)$} the \emph{b-divisor of $T$}\index{b-divisor!b-divisor of a current}.
\end{definition}
This assignment is compatible with push-forwards. Indeed, if $g\colon Z\to Y$ is a modification and $\pi\colon Y\to X$, then
\[
    g_*\Reg\bigl(g^*\pi^*T\bigr)=\Reg(\pi^*T)
\]
by Siu's decomposition and \cref{lma:Siudec_pullpush}. Hence
\[
    g_*\mathbb D(T)_Z=\mathbb D(T)_Y.
\]
Thus $\mathbb D(T)$ is a Weil b-divisor.

We observe that if $T'$ is another closed positive $(1,1)$-current on $X$ and $\lambda\geq 0$, then
\[
\mathbb{D}(T+T')=\mathbb{D}(T)+\mathbb{D}(T'),\quad \mathbb{D}(\lambda T)=\lambda \mathbb{D}(T).
\]
We shall use these identities implicitly in the sequel.

Note that when $T$ has analytic singularities, $\mathbb{D}(T)$ is Cartier.

\begin{theorem}\label{thm:DT}
    Let $T$ be a closed positive $(1,1)$-current on $X$. Then $\mathbb{D}(T)$ is nef. Moreover,
    \begin{equation}\label{eq:volDT}
    \vol T=\vol \mathbb{D}(T).
    \end{equation}
\end{theorem}

\begin{proof}
    Let $\omega$ be a Kähler form on $X$.

    \textbf{Step~1}. Reduce to the case where $T$ is a Kähler current.

    Note that $\mathbb{D}(\omega)$ is the Cartier b-divisor realized by $(X,\{\omega\})$.
    The b-divisors $\mathbb{D}(T+\epsilon\omega)=\mathbb{D}(T)+\epsilon \mathbb{D}(\omega)$ decrease to $\mathbb{D}(T)$ as $\epsilon\searrow 0$. Hence, by \cref{lma:volbdivnet} and \cref{prop:vollinearlimit}, it suffices to prove our assertion with $T+\epsilon\omega$ in place of $T$.

    \textbf{Step~2}. We prove the assertion under the additional assumption that $T$ has analytic singularities.

    Let $\pi\colon Y\rightarrow X$ be a modification so that
    \[
    \pi^*T=[D]+R,
    \]
    where $D$ is an effective $\mathbb{Q}$-divisor on $Y$ and $R$ is a closed positive $(1,1)$-current with locally bounded potentials. See \cref{thm:resolvelogsing} for the existence of $\pi$.
    Then $\mathbb{D}(T)$ is the nef Cartier b-divisor realized by $(\pi,\{R\})$. Note that \eqref{eq:volDT} is immediate in this case.

    \textbf{Step~3}. We prove the assertion for a general Kähler current $T$.
    Next, take a closed smooth real $(1,1)$-form $\theta$ in the cohomology class of $T$ and write $T=\theta_{\varphi}$ for some $\varphi\in \PSH(X,\theta)$. Since $T$ is a Kähler current, we choose the quasi-equisingular approximation $(\varphi_j)_j$ of $\varphi$ so that, after discarding finitely many terms, each $\theta+\ddc\varphi_j$ is a Kähler current. This is possible by Demailly regularization, in the form used in the proof of \cref{thm:charIgoodasclosure}. We claim that
    \begin{equation}\label{eq:conv_bdiv_temp1}
        \mathbb{D}(\theta+\ddc \varphi_j)\to \mathbb{D}(\theta+\ddc \varphi).
    \end{equation}
    By definition of this convergence, we need to establish the following: for every modification $\pi\colon Y\rightarrow X$, we have
    \[
        \left\{\Reg\left(\pi^*(\theta+\ddc\varphi_j)\right)\right\}\to \left\{\Reg\left(\pi^*(\theta+\ddc\varphi)\right)\right\}.
    \]
    This follows immediately from \cref{thm:Lelongcont} if $\Sing (\pi^*T)$ has only finitely many components.

    Fix a Kähler form $\omega_Y$ on $Y$. We want to show that for any $\epsilon>0$, we can find $j_0>0$ so that when $j\geq j_0$,
\begin{equation}\label{eq:Regpivarphij_conv_temp1}
\left\{\Reg\left(\pi^*(\theta+\ddc\varphi_j)\right)\right\}\leq \left\{\Reg\left(\pi^*(\theta+\ddc\varphi)\right)\right\}+\epsilon\{\omega_Y\}.
\end{equation}
Write the divisorial part of $\pi^*\theta+\ddc \pi^*\varphi_j$ and $\pi^*\theta+\ddc \pi^*\varphi$ as
\[
\sum_{i=1}^{\infty}a_i^j [E_i],\quad \sum_{i=1}^{\infty}a_i [E_i].
\]
Then $a_i^j\leq a_i$.

We can find an integer $N>0$ large enough so that, in the pseudo-effective order of cohomology classes,
\[
\sum_{i=N+1}^{\infty}a_i \{E_i\}\leq \frac{\epsilon}{2}\{\omega_Y\}.
\]
By \cref{thm:Lelongcont}, we can take $j_0$ large enough so that for $j>j_0$,
\[
\sum_{i=1}^{N}(a_i-a_i^j)\{E_i\}\leq \frac{\epsilon}{2}\{\omega_Y\}.
\]
Then \eqref{eq:Regpivarphij_conv_temp1} follows, and \eqref{eq:conv_bdiv_temp1} is established.

As a consequence, $\mathbb{D}(T)$ is nef and the volume can be computed using \cref{lma:volbdivnet}:
\[
\vol\mathbb{D}(T)=\lim_{j\to\infty}\vol \mathbb{D}\left( \theta+\ddc\varphi_j \right)=\lim_{j\to\infty}\vol \left(\theta+\ddc\varphi_j\right)=\vol T.
\]
Hence, \eqref{eq:volDT} follows.
\end{proof}


Conversely, we want to realize nef b-divisors as $\mathbb{D}(T)$. We first prove a continuity result.
\begin{proposition}\label{prop:dscontbdiv}
Let $\theta$ be a closed real smooth $(1,1)$-form on $X$, let $(\varphi_i)_{i\in I}$ be a net in $\PSH(X,\theta)$, and let $\varphi\in \PSH(X,\theta)$. If $\varphi_i\xrightarrow{d_S}\varphi$, then
\begin{equation}\label{eq:Dcont1}
    \mathbb{D}\left(\theta_{\varphi_i}\right)\to \mathbb{D}\left(\theta_{\varphi}\right).
\end{equation}
\end{proposition}
\begin{proof}
    Fix a modification $\pi\colon Y\rightarrow X$.
    Fix a Kähler form $\omega$ on $X$.
    It suffices to establish the following:
    \begin{equation}\label{eq:Dcont2}
    \mathbb{D}\left(\theta_{\varphi_i}\right)_Y\to \mathbb{D}\left(\theta_{\varphi}\right)_Y.
    \end{equation}
    As a map from the pseudometric space $\PSH(X,\theta)$ to the finite-dimensional Euclidean space $\mathrm{H}^{1,1}(Y,\mathbb{R})$, the continuity of $\mathbb{D}(\bullet)_Y$ can be tested on sequences. Thus, without loss of generality, we may assume that $(\varphi_i)_i$ is a sequence and $I=\mathbb{Z}_{>0}$. It is enough to prove that every subsequence has a further subsequence for which \eqref{eq:Dcont3} holds.

    By \cref{cor:dSconv_bimero}, $\pi^*\varphi_i\xrightarrow{d_S}\pi^*\varphi$. Hence it suffices to prove the claim on $Y$, so we may assume that $\pi$ is the identity. Therefore, we are reduced to the following assertion:
    \begin{equation}\label{eq:Dcont3}
        \left\{ \Reg \theta_{\varphi_i} \right\}\to \left\{ \Reg \theta_{\varphi} \right\}.
    \end{equation}
    After adding a Kähler form to $\theta$, we may also assume that $\theta_{\varphi}$ is a Kähler current. By \cref{thm:contPI}, replacing all potentials by their $\mathcal{I}$-envelopes does not change the divisorial parts and preserves $d_S$-convergence. The masses are then bounded from below along a tail, so after passing to a subsequence \cref{prop:incanddec} supplies increasing and decreasing envelopes $\eta_i\leq \varphi_i\leq \psi_i$ which both converge to $\varphi$ in $d_S$. Since
    \[
        \mathbb{D}(\theta_{\eta_i})_X\leq \mathbb{D}(\theta_{\varphi_i})_X\leq \mathbb{D}(\theta_{\psi_i})_X,
    \]
    it is enough to treat the case where $(\varphi_i)_i$ is monotone.

    Write the divisorial parts of $\theta_{\varphi_i}$ and $\theta_{\varphi}$ as
    \[
        \sum_{m=1}^{\infty}a_m^i[D_m],\quad \sum_{m=1}^{\infty}a_m[D_m].
    \]
    By \cref{thm:Lelongcont}, $a_m^i\to a_m$ for each fixed $m$. If $(\varphi_i)_i$ is decreasing, then $a_m^i\leq a_m$ for all $i,m$. The finite-tail argument in the proof of \eqref{eq:Regpivarphij_conv_temp1} then gives
    \[
        \left\{\Reg\theta_{\varphi}\right\}\leq \left\{\Reg\theta_{\varphi_i}\right\}\leq \left\{\Reg\theta_{\varphi}\right\}+\epsilon\{\omega\}
    \]
    for all sufficiently large $i$.

    If $(\varphi_i)_i$ is increasing, then $a_m\leq a_m^i\leq a_m^1$. We first choose $N$ so that the tails of both $\sum_m a_m[D_m]$ and $\sum_m a_m^1[D_m]$ are at most $(\epsilon/3)\{\omega\}$ in the pseudo-effective order. The convergence of the finitely many coefficients $a_m^i$, $m\leq N$, then gives
    \[
        \left\{\Reg\theta_{\varphi_i}\right\}\leq \left\{\Reg\theta_{\varphi}\right\}\leq \left\{\Reg\theta_{\varphi_i}\right\}+\epsilon\{\omega\}
    \]
    for all sufficiently large $i$. This proves \eqref{eq:Dcont3}.
\end{proof}

\begin{lemma}\label{lma:pushIpo}
    Let $\pi\colon X\rightarrow Z$ be a proper bimeromorphic morphism from $X$ to a Kähler manifold $Z$. Consider non-divisorial closed positive $(1,1)$ currents $T,S$ on $X$ in the same cohomology class. If $T\preceq_{\mathcal{I}} S$, then $\pi_*T\preceq_{\mathcal{I}} \pi_*S$.
\end{lemma}
\begin{proof}
    By Hironaka's Chow lemma \cref{thm:HironakaChow} and \cref{lma:Ipartorderinv}, we may assume that $\pi$ is a modification.

    By \cref{lma:Siudec_pullpush},
    \[
    \pi^*\pi_*T=T+F_T,\quad \pi^*\pi_*S=S+F_S,
    \]
    where $F_T$ and $F_S$ are effective $\pi$-exceptional divisors. Since $\{T\}=\{S\}$, the exceptional divisors $F_T$ and $F_S$ have the same cohomology class. Factoring $\pi$ into blow-ups and using \cref{prop:blowup_cone} at each step, the classes of $\pi$-exceptional prime divisors are linearly independent in the relative cohomology; hence $F_T=F_S$. It follows that
    \[
    \pi^*\pi_*T=T+F_T\preceq_{\mathcal{I}}S+F_S=\pi^*\pi_*S.
    \]
    Our assertion then follows from \cref{lma:Ipartorderinv}.
\end{proof}
\begin{theorem}\label{thm:bigandnefreal}
    Each big and nef b-divisor $\mathbb{D}$ over $X$ can be realized as $\mathbb{D}(T)$ for some $T\in \mathbb{D}_X$. Furthermore, we may always assume that $T$ is $\mathcal{I}$-good.
\end{theorem}
Note that $T$ is not unique. The current $T$ is necessarily non-divisorial.

\begin{proof}
    Fix a big and nef b-divisor $\mathbb{D}$ over $X$.

    For each $\pi\colon Y\rightarrow X$, we take a current with minimal singularities $T_Y$ in $\mathbb{D}_Y$.
    We claim that $\mathbb{D}(\pi_*T_Y)$ coincides with $\mathbb{D}$ up to the level of $Y$: For any modification $\pi'\colon Z\rightarrow X$ dominated by $\pi$ through a morphism $g\colon Y\rightarrow Z$, we have
    \[
    \mathbb{D}_Z=\mathbb{D}(\pi_*T_Y)_Z.
    \]
    The notations are summarized in the following commutative diagram:
    \begin{equation}\label{eq:domi2}
    \begin{tikzcd}
Y \arrow[rr,"g"] \arrow[rd, "\pi"'] &   & Z \arrow[ld, "\pi'"] \\
                                 & X. &
\end{tikzcd}
    \end{equation}
    After unfolding the definitions, this means
    \[
        \Reg(\pi'^*\pi_*T_Y)\in \mathbb{D}_Z.
    \]
    Note that
    \[
        \Reg(\pi'^*\pi_*T_Y)=\Reg(\pi'^* \pi'_* g_*T_Y).
    \]
    Due to \cref{prop:cones_modif} and \cref{prop:cones_modif2}, we know that $\mathbb{D}_Y$ is modified nef and big. In particular, $T_Y$ is non-divisorial, hence so is $g_*T_Y$ by \cref{lma:ndvpush}. It follows from \cref{lma:Siudec_pullpush} that
    \[
    \Reg(\pi'^* \pi'_* g_*T_Y)=\Reg(g_*T_Y)=g_*T_Y\in \mathbb{D}_Z.
    \]
    Note that
    \begin{equation}\label{eq:vol_unif_bound}
    \vol T_Y\geq \vol \mathbb{D}>0.
    \end{equation}

    Next we claim that the $P$-singularity types of the net $(\pi_*T_Y)_Y$ form a decreasing net.

    To see this, let us fix a diagram as \eqref{eq:domi2}. We need to show that
    \[
    \pi_*T_Y\preceq_P \pi'_*T_Z.
    \]
    Since $T_Z$ has minimal singularities, we have $g_*T_Y\preceq_{\mathcal{I}} T_Z$. In particular, \cref{lma:pushIpo} guarantees that $\pi_*T_Y\preceq_{\mathcal{I}} \pi'_*T_Z$.
    But by \cref{cor:Igoodpush} and \cref{ex:ImodelIgood}, both $\pi_*T_Y$ and $\pi'_*T_Z$ are $\mathcal{I}$-good, so there is no difference between the $P$-preorder and the $\mathcal{I}$-preorder in this case. Our assertion follows.

    Next observe that the net $(\pi_*T_Y)_Y$ is decreasing with respect to the $P$-preorder and has a $d_S$-limit as a consequence of \eqref{eq:vol_unif_bound} and \cref{cor:completenessdS}.
    Take a closed positive $(1,1)$-current $T\in \mathbb{D}_X$ such that
    \[
    \pi_*T_Y\xrightarrow{d_S}T.
    \]
    It follows from \cref{prop:dscontbdiv} that
    \[
    \mathbb{D}(\pi_*T_Y)\to \mathbb{D}(T).
    \]
    Therefore, we conclude that
    \[
    \mathbb{D}(T)=\mathbb{D}.
    \]
    Thanks to \cref{thm:DT}, $\vol T>0$. Write $T=\theta+\ddc \varphi$ for some $\varphi\in \PSH(X,\theta)$, then
    \[
    T'\coloneqq \theta+\ddc P_{\theta}[\varphi]_{\mathcal{I}}
    \]
    is $\mathcal{I}$-good, non-divisorial and $\mathbb{D}(T')=\mathbb{D}(T)$.
\end{proof}

\begin{corollary}\label{cor:mn_level_nefb}
    Let $\mathbb{D}$ be a b-divisor over $X$. Then the following are equivalent:
    \begin{enumerate}
        \item\label{itm:cor:mn_level_nefb:1} $\mathbb{D}$ is nef;
        \item\label{itm:cor:mn_level_nefb:2} for each modification $\pi\colon Y\rightarrow X$, the class $\mathbb{D}_Y$ is modified nef.
    \end{enumerate}
\end{corollary}
\begin{proof}
    \ref{itm:cor:mn_level_nefb:1} $\implies$ \ref{itm:cor:mn_level_nefb:2}. Suppose that \ref{itm:cor:mn_level_nefb:1} holds. Fix a modification $\pi\colon Y\rightarrow X$. Then we need to show that $\mathbb{D}_Y$ is modified nef. Since modified nefness is a closed condition, after approximating $\mathbb{D}$ by nef Cartier b-divisors, we may assume that $\mathbb{D}$ itself is a nef Cartier b-divisor. We can then find a modification $\pi'\colon Z\rightarrow X$ dominating $\pi$ so that $\mathbb{D}$ is realized by a nef class $\alpha$ on $Z$. Then $\mathbb{D}_Y$ is nothing but the pushforward of $\alpha$, and hence modified nef thanks to \cref{prop:cones_modif}.

    \ref{itm:cor:mn_level_nefb:2} $\implies$ \ref{itm:cor:mn_level_nefb:1}. Suppose that \ref{itm:cor:mn_level_nefb:2} holds. We need to show that $\mathbb{D}$ is nef. Fix a Kähler form $\omega$ on $X$. It suffices to show that $\mathbb{D}+\epsilon\mathbb{D}(\omega)$ is nef for each $\epsilon>0$. After replacing $\mathbb{D}$ by the latter, we may further assume that $\mathbb{D}_Y$ is big for each modification $\pi\colon Y\rightarrow X$, and $\vol \mathbb{D}_Y$ has a uniform positive lower bound.
    In this case, the argument of \cref{thm:bigandnefreal} shows that $\mathbb{D}=\mathbb{D}(T)$ for some closed positive $(1,1)$-current in $\mathbb{D}_X$. Therefore, $\mathbb{D}$ is nef, thanks to \cref{thm:DT}.
\end{proof}


Let $\alpha$ be a modified nef class on $X$. We write $\mathcal{G}(\alpha)$ for the set of closed positive $(1,1)$-currents $T$ on $X$ with $T=\Reg T\in \alpha$ and $\vol T>0$.
\begin{theorem}\label{thm:bdiv_current_corr}
    There is a natural bijection from $\mathcal{G}(\alpha)/\sim_{\mathcal{I}}$ to the set of big and nef b-divisors $\mathbb{D}$ over $X$ with $\mathbb{D}_X=\alpha$.
\end{theorem}
\begin{proof}


    Given $T\in \mathcal{G}(\alpha)$, we associate the b-divisor $\mathbb{D}(T)$. It is big and nef due to \cref{thm:DT}. This map clearly descends to $\mathcal{G}(\alpha)/\sim_{\mathcal{I}}$.

    This map is surjective by \cref{thm:bigandnefreal}. Now we show that it is injective. Let $T,T'\in \mathcal{G}(\alpha)$. Assume that $\mathbb{D}(T)=\mathbb{D}(T')$, we want to show that $T\sim_{\mathcal{I}} T'$.

    Let $E$ be a prime divisor over $X$, it suffices to show that
    \begin{equation}\label{eq:nuTE}
    \nu(T,E)=\nu(T',E).
    \end{equation}
    We may assume that $E$ is not a prime divisor on $X$, as otherwise both sides vanish.

    Choose a sequence of blow-ups with smooth \emph{connected} centers
    \[
    Y\coloneqq X_k\rightarrow X_{k-1}\rightarrow \cdots \rightarrow X_0\coloneqq X
    \]
    so that $E$ is a prime divisor on $Y$, exceptional with respect to $X_k\rightarrow X_{k-1}$. Denote the composition by $\pi\colon Y\rightarrow X$.
    Thanks to \cref{prop:blowup_cone},
    \[
    \mathrm{H}^{1,1}(X_k,\mathbb{R})=\mathrm{H}^{1,1}(X_{k-1},\mathbb{R})\oplus \mathbb{R}\{E_k\},
    \]
    where $E_k=E$ is the exceptional divisor of $X_k\rightarrow X_{k-1}$.

    By induction,
    \[
    \mathrm{H}^{1,1}(Y,\mathbb{R})=\mathrm{H}^{1,1}(X,\mathbb{R})\oplus \bigoplus_{i=1}^{k} \mathbb{R} \{E_i\},
    \]
    where $E_i$ is the exceptional divisor of $X_i\rightarrow X_{i-1}$. Now by \cref{lma:Siudec_pullpush},
    \begin{equation}\label{eq:RegpistarT}
    \Reg \pi^*T=\pi^*T-\sum_{i=1}^k \nu(T,E_i)[E_i].
    \end{equation}
    In particular, the cohomology class of $\Reg \pi^*T$ determines $\nu(T,E)$. Hence, \eqref{eq:nuTE} follows.
\end{proof}


\begin{corollary}\label{cor:nef_b_degenerate_1}
    The set of nef b-divisors over $X$ with root $\alpha$ can be naturally identified with
    \[
    \varprojlim_{\omega} \left( \mathcal{G}(\alpha+\omega)/\sim_{\mathcal{I}} \right),
    \]
    where $\omega$ runs over the directed set of Kähler forms on $X$ (with respect to the preorder of reverse domination), and given two Kähler forms $\omega\leq \omega'$ the transition map
    \[
    \mathcal{G}(\alpha+\omega)/\sim_{\mathcal{I}}\rightarrow \mathcal{G}(\alpha+\omega')/\sim_{\mathcal{I}}
    \]
    is induced by the map $\mathcal{G}(\alpha+\omega)\rightarrow \mathcal{G}(\alpha+\omega')$ sending $T$ to $T+\omega'-\omega$.
\end{corollary}
\begin{proof}
For each Kähler form $\omega$, \cref{thm:bdiv_current_corr} identifies $\mathcal G(\alpha+\omega)/\sim_{\mathcal I}$ with the set of big and nef b-divisors whose root on $X$ is $\alpha+\{\omega\}$. The transition map is exactly addition of $\mathbb D(\omega'-\omega)$. Passing to the inverse limit  gives precisely nef b-divisors with root $\alpha$.
\end{proof}

It is tempting to extend these results to more general currents, not necessarily non-divisorial ones. For this purpose, we introduce the following definition:
\begin{definition}\label{def:augmented_nefb}
    An \emph{augmented nef b-divisor}\index{augmented nef b-divisor} over $X$ is a pair $(\mathbb{D},D)$, where
    \begin{enumerate}
        \item\label{itm:def:augmented_nefb:1} $\mathbb{D}$ is a nef b-divisor over $X$;
        \item\label{itm:def:augmented_nefb:2} $D$ is a formal sum
        \[
        D=\sum_{E} c_E E,\quad c_E\in \mathbb{R}_{\geq 0},
        \]
        where $E$ runs over the set of prime divisors on $X$,
    \end{enumerate}
    such that the following condition is satisfied:
    \[
    \sum_{E}c_E \{E\}
    \]
    is convergent as a sum in $\mathrm{H}^{1,1}(X,\mathbb{R})$.

    The \emph{cohomology class} of $(\mathbb{D},D)$ is defined as
    \[
    \mathbb{D}_X+\sum_{E}c_E \{E\}\in \mathrm{H}^{1,1}(X,\mathbb{R}).
    \]
    The \emph{volume} of $(\mathbb{D},D)$ is defined as
    \[
    \vol (\mathbb{D},D)=\vol \mathbb{D}.
    \]
\end{definition}

Recall that we defined $\mathcal{Z}_+(X,\alpha)$ as the set of closed positive $(1,1)$-currents $T\in \alpha$ in \cref{def:spaceofcurrents}. We introduce a further notation here:
\[
\mathcal{Z}_+(X,\alpha)_{>0}\coloneqq \left\{T\in \mathcal{Z}_+(X,\alpha): \vol T>0\right\}.
\]
\index{$\mathcal{Z}_+(X,\alpha)_{>0}$}

\begin{corollary}\label{cor:Iequiv_curr_b_div}
    There is a canonical bijection between the following sets:
    \begin{enumerate}
        \item\label{itm:cor:Iequiv_curr_b_div:1} The set $\mathcal{Z}_+(X,\alpha)_{>0}/\sim_{\mathcal{I}}$;
        \item\label{itm:cor:Iequiv_curr_b_div:2} the set of augmented nef b-divisors over $X$ with positive volume with cohomology class $\alpha$.
    \end{enumerate}
\end{corollary}
More precisely, given a current $T\in \mathcal{Z}_+(X,\alpha)_{>0}$, we associate $(\mathbb{D},D)$ as follows:
\[
\mathbb{D}=\mathbb{D}(T),\quad D=\sum_E \nu(T,E)E,
\]
where $E$ runs over the set of prime divisors on $X$.
\begin{proof}
    This is a direct consequence of \cref{thm:bdiv_current_corr} and Siu's decomposition \cref{lma:Siudec}.
\end{proof}

In fact, \cref{cor:Iequiv_curr_b_div} can also be reformulated in an elementary manner, without referring to b-divisors at all. Suppose that we are given a big cohomology class $\alpha$ on $X$ and a non-negative real number $c_E$ for each prime divisor $E$ over $X$. A natural question is: When is there a closed positive $(1,1)$-current $T\in \alpha$ with positive volume such that $\nu(T,E)=c_E$ for every $E$? Then \cref{cor:Iequiv_curr_b_div} says that $T$ exists if and only if the following two conditions hold:
\begin{enumerate}
    \item\label{itm:chapter-bdiv:L472:1}
    \[
    \lim_{\pi\colon Y\rightarrow X} \vol\left( \pi^*\alpha-\sum_{E\subseteq Y} c_E \{E\} \right)>0\footnotemark,
    \]
    \footnotetext{In particular, implicitly, the sum $\sum_{E\subseteq Y} c_E \{E\}$ converges in $\mathrm{H}^{1,1}(Y,\mathbb{R})$.}
    where $\pi$ runs over all modifications of $X$;
    \item\label{itm:chapter-bdiv:L472:2} for each modification $\pi\colon Y\rightarrow X$, the class $ \pi^*\alpha-\sum_{E\subseteq Y} c_E \{E\}$ is modified nef.
\end{enumerate}




Similarly, we have the following generalization of \cref{cor:nef_b_degenerate_1}.
\begin{corollary}\label{cor:nef_b_degenerate_2}
    There is a canonical bijection between the following two sets:
    \begin{enumerate}
        \item\label{itm:cor:nef_b_degenerate_2:1} The set of
        \[
            \varprojlim_{\omega} \left( \mathcal{Z}_+(X,\alpha+\omega)_{>0}/\sim_{\mathcal{I}} \right),
        \]
        where $\omega$ runs over the directed set of Kähler forms on $X$;
    \item\label{itm:cor:nef_b_degenerate_2:2} the set of augmented nef b-divisors over $X$ with cohomology class $\alpha$.
    \end{enumerate}
\end{corollary}
\begin{proof}
Apply \cref{cor:nef_b_degenerate_1} to the non-divisorial parts and use Siu's decomposition on $X$ to record the divisorial part. The convergence condition in \cref{def:augmented_nefb} is exactly the condition that the roots of the non-divisorial parts converge to the prescribed cohomology class $\alpha$.
\end{proof}

\begin{corollary}\label{cor:nefbCequiv}
    Let $\mathbb{D}$ be a Cartier b-divisor over $X$. Then $\mathbb{D}$ is nef in the sense of \cref{def:nefCar} if and only if it is nef in the sense of \cref{def:nefb}.
\end{corollary}

\begin{proof}
    We only handle the nontrivial implication. Assume that $\mathbb{D}$ is nef in the sense of \cref{def:nefb}. We want to show that $\mathbb{D}$ is nef in the sense of \cref{def:nefCar}. After replacing $X$ by a realization of $\mathbb{D}$, we may assume that $\mathbb{D}$ is realized by $(X,\alpha)$ for some cohomology class $\alpha\in \mathrm{H}^{1,1}(X,\mathbb{R})$.

    We first assume that $\mathbb{D}$ is big. Choose a non-divisorial closed positive $(1,1)$-current $T$ on $X$ such that $\mathbb{D}=\mathbb{D}(T)$. Now $\mathbb{D}=\mathbb{D}(T)$ means that for each modification $\pi\colon Y\rightarrow X$, the current $\pi^*T$ is non-divisorial. In particular, $T$ has vanishing generic Lelong number along each prime divisor over $X$, see \eqref{eq:RegpistarT}. Equivalently, $T$ has vanishing Lelong number everywhere. By Demailly's regularization theorem, a pseudo-effective class carrying a positive current with vanishing Lelong number at every point is nef. Hence $\alpha=\{T\}$ is nef.

    In the general case, fix a Kähler form $\omega$ on the realization of $\mathbb{D}$. For every $\epsilon>0$, the Cartier b-divisor $\mathbb{D}+\epsilon\mathbb{D}(\omega)$ is big and nef in the sense of \cref{def:nefb}, hence is nef in the sense of \cref{def:nefCar} by the big case. Thus $\alpha+\epsilon\{\omega\}$ is nef for all $\epsilon>0$, and the closedness of the nef cone shows that $\alpha$ is nef.
\end{proof}

\begin{corollary}\label{cor:Dorder_Iorder}
    Let $T$ and $T'$ be non-divisorial closed positive $(1,1)$-currents on $X$. Suppose that $\{T\}=\{T'\}$. Then the following are equivalent:
    \begin{enumerate}
        \item\label{itm:cor:Dorder_Iorder:1} $\mathbb{D}(T)\leq \mathbb{D}(T')$;
        \item\label{itm:cor:Dorder_Iorder:2} $T\preceq_{\mathcal{I}}T'$.
    \end{enumerate}
\end{corollary}
\begin{proof}
    Let $E$ be a prime divisor over $X$. Choose a sequence of blow-ups with smooth connected centers as in the proof of \cref{thm:bdiv_current_corr}, so that $E$ appears as one of the exceptional prime divisors on the resulting model $\pi\colon Y\rightarrow X$. By \eqref{eq:RegpistarT}, the difference
    \[
        \mathbb{D}(T')_Y-\mathbb{D}(T)_Y
    \]
    is represented in the relative cohomology by
    \[
        \sum_i\left(\nu(T,E_i)-\nu(T',E_i)\right)\{E_i\}.
    \]
    Using \cref{prop:blowup_cone} inductively, or equivalently the negativity lemma for exceptional classes, pseudo-effectivity of this exceptional class is equivalent to the non-negativity of all coefficients. Hence $\mathbb{D}(T)\leq \mathbb{D}(T')$ is equivalent to
    \[
        \nu(T,E)\geq \nu(T',E)
    \]
    for every prime divisor $E$ over $X$, which is exactly $T\preceq_{\mathcal{I}}T'$ by \cref{prop:Iequivchar2}.
\end{proof}

\begin{corollary}\label{cor:nef_Cartierseq}
    Let $\mathbb{D}$ be a nef b-divisor over $X$. Then there is a decreasing sequence of nef and big Cartier b-divisors $\mathbb{D}_i$ over $X$ with limit $\mathbb{D}$.
\end{corollary}
\begin{proof}
    Take a Kähler form $\omega$ on $X$. By \cref{thm:bigandnefreal}, for each $i>0$, we can find a non-divisorial Kähler current $T_i\in \mathbb{D}_X+i^{-1}\{\omega\}$ such that
    \[
    \mathbb{D}(T_i)=\mathbb{D}+i^{-1}\mathbb{D}(\omega).
    \]
    We observe that $T_i$ and $T_{i+1}$ have the same $\mathcal{I}$-singularity type. Indeed, applying \cref{cor:Dorder_Iorder} to $T_i$ and $T_{i+1}+(i^{-1}-(i+1)^{-1})\omega$ shows that these two currents are $\mathcal{I}$-equivalent; adding the smooth form does not change the singularity type.
    Let $(T_i^j)_j$ be quasi-equisingular approximations of $T_i$ such that
    \begin{enumerate}
        \item\label{itm:chapter-bdiv:L544:1}
        $T_i^j$ is a Kähler current in $\mathbb{D}_X+i^{-1}\{\omega\}$ for $j\geq j_0(i)$, and
        \item\label{itm:chapter-bdiv:L544:2} for each fixed $j$, the singularity type of $T_i^j$ is independent of $i$.
    \end{enumerate}
    Note that \ref{itm:chapter-bdiv:L544:2} is possible by using the Bergman kernel construction of the quasi-equisingular approximations, since the $T_i$'s have the same multiplier ideals.

    For fixed $j$, the b-divisor
    \[
        \mathbb{B}_j\coloneqq \mathbb{D}(T_i^j)-i^{-1}\mathbb{D}(\omega)
    \]
    is independent of $i$. Since the quasi-equisingular approximations are decreasing in $j$, the sequence $(\mathbb{B}_j)_j$ is decreasing, and \cref{prop:dscontbdiv} gives $\mathbb{B}_j\to \mathbb{D}$. It suffices to take $\mathbb{D}_i=\mathbb{B}_{j_i}+i^{-1}\mathbb{D}(\omega)=\mathbb{D}(T_i^{j_i})$, where $j_i$ is chosen to be a sufficiently fast increasing sequence of positive integers with $j_i\geq j_0(i)$.
\end{proof}


\section{Properties of the intersection product}\label{sec:int_b_div}
Let $X$ be a connected compact Kähler manifold of dimension $n$. Recall that the mixed volume of currents is studied in \cref{sec:mix_vol}.

\begin{definition}\label{def:DF_ext2}
    Let $\mathbb{D}_1,\ldots,\mathbb{D}_n$ be big and nef b-divisors over $X$. Then we define their intersection as
    \index{b-divisor!intersection product}
    \[
    \left( \mathbb{D}_1,\ldots,\mathbb{D}_n \right)\coloneqq \vol (T_1,\ldots,T_n),
    \]
    \index{$\left( \mathbb{D}_1,\ldots,\mathbb{D}_n \right)$}%
    where $T_1,\ldots,T_n$ are closed positive $(1,1)$-currents in $\mathbb{D}_{1,X},\ldots,\mathbb{D}_{n,X}$ respectively such that $\mathbb{D}(T_i)=\mathbb{D}_{i}$.
\end{definition}
The definition makes sense thanks to \cref{thm:bigandnefreal}. It does not depend on the choices of $T_1,\ldots,T_n$ since they are uniquely defined up to $\mathcal{I}$-equivalence, as proved in \cref{thm:bdiv_current_corr}.



\begin{lemma}\label{lma:DFlinear}\leavevmode
    \begin{enumerate}
        \item\label{itm:lma:DFlinear:1} The product in \cref{def:DF_ext2} is symmetric in its $n$ variables.
        \item\label{itm:lma:DFlinear:2} Let $\mathbb{D}_1,\ldots,\mathbb{D}_n,\mathbb{D}'_1$ be big and nef b-divisors over $X$. Then
        \[
        \left(\mathbb{D}_1+\mathbb{D}_1',\ldots,\mathbb{D}_n \right)=\left(\mathbb{D}_1,\ldots,\mathbb{D}_n \right)+\left(\mathbb{D}_1',\ldots,\mathbb{D}_n \right).
        \]
        \item\label{itm:lma:DFlinear:3} Let $\mathbb{D}_1,\ldots,\mathbb{D}_n$ be big and nef b-divisors over $X$ and $\lambda\geq 0$. Then
        \[
        \left(\lambda \mathbb{D}_1,\ldots,\mathbb{D}_n \right)=\lambda \left(\mathbb{D}_1,\ldots,\mathbb{D}_n \right).
        \]
        \item\label{itm:lma:DFlinear:4} Let $\mathbb{D}_1,\ldots,\mathbb{D}_n$ and $\mathbb{D}'$ be big and nef b-divisors over $X$ such that $\mathbb{D}_1\leq \mathbb{D}'$. Then
        \[
            \left(\mathbb{D}_1,\ldots,\mathbb{D}_n  \right)\leq \left(\mathbb{D}',\mathbb{D}_2,\ldots,\mathbb{D}_n  \right).
        \]
    \end{enumerate}
\end{lemma}
\begin{proof}
    \ref{itm:lma:DFlinear:1}, \ref{itm:lma:DFlinear:2} and \ref{itm:lma:DFlinear:3} follow immediately from the corresponding results for mixed volumes, as proved in \cref{prop:vollinearlimit}.

    \ref{itm:lma:DFlinear:4} Take $\mathcal{I}$-good non-divisorial closed positive $(1,1)$-currents $T_1,\ldots,T_n$ and $T'$ so that $\mathbb{D}(T_i)=\mathbb{D}_i$ for all $i=1,\ldots,n$ and $\mathbb{D}(T')=\mathbb{D}'$.
    Furthermore, we may assume that the $T_i$'s and $T'$ are Kähler currents by the perturbation argument.

    Let $(T_{i}^j)_j$ be a quasi-equisingular approximation of $T_i$ for $i=2,\ldots,n$. It follows from \cref{thm:convdS} that
    \[
    \int_X T_1\wedge \cdots \wedge T_n=\lim_{j\to\infty} \int_X T_1\wedge T_2^j\wedge\cdots \wedge T_n^j.
    \]
    It suffices to show that for all $j\geq 1$,
    \[
    \int_X T_1\wedge T_2^j\wedge\cdots \wedge T_n^j\leq \int_X T'\wedge T_2^j\wedge\cdots \wedge T_n^j.
    \]
    Therefore, we have reduced to the case where $T_2,\ldots,T_n$ have analytic singularities. After a resolution, we may assume that they have log singularities along $\mathbb{Q}$-divisors. By \itemref{itm:prop:vollinearlimit:5}, we can further reduce to the case where $T_2,\ldots,T_n$ have bounded local potentials. Perturbing $T_2,\ldots,T_n$ by a Kähler form, we may further assume that $\{T_2\},\ldots, \{T_n\}$ are Kähler classes. By \itemref{itm:prop:vollinearlimit:4}, we finally reduce to the case where $T_2,\ldots,T_n$ are Kähler forms. In this case, the assertion follows from the pseudo-effectivity of $\{T'\}-\{T_1\}$.
\end{proof}


In general, we can extend the definition of the intersection product to nef b-divisors as follows:
\begin{definition}\label{def:DF_ext}
    Let $\mathbb{D}_1,\ldots,\mathbb{D}_n$ be nef b-divisors over $X$. Then we define their intersection as
    \index{b-divisor!intersection product}
    \begin{equation}\label{eq:mixed_product_limit_def}
        \left( \mathbb{D}_1,\ldots,\mathbb{D}_n \right)\coloneqq \lim_{\epsilon \to 0+}\Bigl( \mathbb{D}_1+\epsilon \mathbb{D}(\omega),\ldots,\mathbb{D}_n+\epsilon \mathbb{D}(\omega) \Bigr),
    \end{equation}
    \index{$\left( \mathbb{D}_1,\ldots,\mathbb{D}_n \right)$}%
    where $\omega$ is a Kähler form on $X$.

    Note that as $\epsilon\to 0+$, $\left( \mathbb{D}_1+\epsilon \mathbb{D}(\omega),\ldots,\mathbb{D}_n+\epsilon \mathbb{D}(\omega) \right)$ is decreasing, as a consequence of \itemref{itm:lma:DFlinear:4}. Similarly, by \itemref{itm:lma:DFlinear:4} again, the definition is independent of the choice of $\omega$.
\end{definition}


\begin{proposition}\label{prop:DFlinear}\leavevmode
    \begin{enumerate}
        \item\label{itm:prop:DFlinear:1} The product in \cref{def:DF_ext} is symmetric in its $n$ variables.
        \item\label{itm:prop:DFlinear:2} Let $\mathbb{D}_1,\ldots,\mathbb{D}_n,\mathbb{D}'_1$ be nef b-divisors over $X$. Then
        \[
        \left(\mathbb{D}_1+\mathbb{D}_1',\ldots,\mathbb{D}_n \right)=\left(\mathbb{D}_1,\ldots,\mathbb{D}_n \right)+\left(\mathbb{D}_1',\ldots,\mathbb{D}_n \right).
        \]
        \item\label{itm:prop:DFlinear:3} Let $\mathbb{D}_1,\ldots,\mathbb{D}_n$ be nef b-divisors over $X$ and $\lambda\geq 0$. Then
        \[
        \left(\lambda \mathbb{D}_1,\ldots,\mathbb{D}_n \right)=\lambda \left(\mathbb{D}_1,\ldots,\mathbb{D}_n \right).
        \]
        \item\label{itm:prop:DFlinear:4} Let $\mathbb{D}_1,\ldots,\mathbb{D}_n$ and $\mathbb{D}'$ be nef b-divisors over $X$ such that $\mathbb{D}_1\leq \mathbb{D}'$. Then
        \[
            \left(\mathbb{D}_1,\ldots,\mathbb{D}_n  \right)\leq \left(\mathbb{D}',\mathbb{D}_2,\ldots,\mathbb{D}_n  \right).
        \]
        \item\label{itm:prop:DFlinear:5} Let $T_1,\ldots,T_n$ be closed positive $(1,1)$-currents on $X$. Then
        \[
            \Bigl( \mathbb{D}(T_1),\ldots,\mathbb{D}(T_n) \Bigr)= \vol (T_1,\ldots,T_n).
        \]
    \end{enumerate}
\end{proposition}
\begin{proof}
    \ref{itm:prop:DFlinear:1}, \ref{itm:prop:DFlinear:2}, \ref{itm:prop:DFlinear:3} and \ref{itm:prop:DFlinear:4} These are immediate consequences of the corresponding results in the big and nef case established in \cref{lma:DFlinear}.

    \ref{itm:prop:DFlinear:5} By \eqref{eq:volmixedgeneral} and \eqref{eq:mixed_product_limit_def}, we may assume that $T_1,\ldots,T_n$ are Kähler currents. In this case, the assertion follows from the definitions.
\end{proof}


\begin{lemma}\label{lma:perturbDF}
    Let $\omega$ be a Kähler form on $X$. Fix a compact set $K\subseteq \mathrm{H}^{1,1}(X,\mathbb{R})$.
    Let $\mathbb{D}_1,\ldots,\mathbb{D}_n$ be nef b-divisors over $X$ such that $\mathbb{D}_{i,X}\in K$ for each $i=1,\ldots,n$. Then there is a constant $C$ depending only on $X,K,\{\omega\}$ such that for any $\epsilon\in [0,1]$, we have
    \[
    0\leq \Bigl( \mathbb{D}_1+\epsilon \mathbb{D}(\omega),\ldots,\mathbb{D}_n+\epsilon \mathbb{D}(\omega) \Bigr)-\left( \mathbb{D}_1,\ldots,\mathbb{D}_n \right)\leq C\epsilon.
    \]
\end{lemma}
\begin{proof}
    By \itemref{itm:lma:DFlinear:4}, the left-hand side is non-negative. For the upper bound, by the multi-linearity of the product proved in \cref{prop:DFlinear}, we can rewrite the difference as a sum of $2^n-1$ terms like
    \[
        \epsilon^{m} \Bigl( \mathbb{D}(\omega)^{m}, \mathbb{D}_{i_1},\ldots, \mathbb{D}_{i_{n-m}} \Bigr),
    \]
    where $J\subseteq\{1,\ldots,n\}$ is a non-empty subset of cardinality $m$, and $\{i_1,\ldots,i_{n-m}\}$ is the complement of $J$. We need to show that these coefficients are uniformly bounded by a constant depending only on the given data. For this purpose, we add $\mathbb{D}(\omega)$ to each $\mathbb{D}_i$ and then replace $K$ by $K+\{\omega\}$. Therefore, we reduce to the case where each $\mathbb{D}_i$ is nef and big, and hence $\mathbb{D}_i=\mathbb{D}(T_i)$ for some $\mathcal{I}$-good closed positive $(1,1)$-currents on $X$ with positive volumes. Then we need to bound
    \[
        \int_X \omega^m\wedge T_{i_1}\wedge \cdots \wedge T_{i_{n-m}}.
    \]
    For this purpose, let $D\geq 0$ be the minimal non-negative number so that $D\{\omega\}-\alpha$ is nef for any $\alpha\in K$. Then the integral is bounded by the cohomological intersection $\{\omega\}^m\cap (D\{\omega\})^{n-m}$. Our assertion follows.
\end{proof}




We first make a consistency check.
\begin{proposition}\label{prop:bdiv:L678}
    Suppose that $\mathbb{D}$ is a nef b-divisor over $X$. Then
    \[
    \left( \mathbb{D},\ldots,\mathbb{D} \right)=\vol \mathbb{D}.
    \]
\end{proposition}
\begin{proof}
    Using \cref{lma:perturbDF} and \cref{lma:volbdivnet}, we may easily reduce to the case where $\mathbb{D}$ is nef and big. In this case, take a non-divisorial closed positive $(1,1)$-current $T$ in $\mathbb{D}_X$ such that $\mathbb{D}(T)=\mathbb{D}$. Then we need to show that
    \[
    \vol \mathbb{D}=\vol T,
    \]
    which is proved in \cref{thm:DT}.
\end{proof}

\begin{proposition}\label{prop:bdiv:L692}
    Let $\mathbb{D}_1,\ldots,\mathbb{D}_n$ be nef b-divisors over $X$. Then
    \[
    \left(\mathbb{D}_1,\ldots,\mathbb{D}_n \right)\geq \prod_{i=1}^n \left( \vol \mathbb{D}_i\right)^{1/n}.
    \]
\end{proposition}
\begin{proof}
    We may assume that $\vol \mathbb{D}_i>0$ for each $i=1,\ldots,n$ since there is nothing to prove otherwise. In this case, our assertion follows from \cref{prop:logvol}.
\end{proof}

\begin{proposition}\label{prop:DFusc}
    The product in \cref{def:DF_ext} is upper semicontinuous in the following sense. Suppose that $(\mathbb{D}_i^j)_{j\in J}$ are nets of nef b-divisors over $X$ with limits $\mathbb{D}_i$ for each $i=1,\ldots,n$. Then
    \begin{equation}\label{eq:limsupint_prod}
    \varlimsup_{j\in J}\left(\mathbb{D}_1^j,\ldots,\mathbb{D}_n^j \right)\leq  \left( \mathbb{D}_1,\ldots,\mathbb{D}_n \right).
    \end{equation}
\end{proposition}
\begin{proof}
    \textbf{Step~1}. We first assume that the $\mathbb D_i^j$'s and the $\mathbb D_i$'s are all big.

    Choose $\mathcal I$-good non-divisorial currents $T_i^j$ and $T_i$ such that $\mathbb D(T_i^j)=\mathbb D_i^j$ and $\mathbb D(T_i)=\mathbb D_i$. Fix a Kähler form $\omega$ on $X$ and fix $\epsilon>0$. Since $\mathbb D_{i,X}^j\to\mathbb D_{i,X}$, for all sufficiently large $j$ the class
    \[
        \mathbb D_{i,X}+\epsilon\{\omega\}-\mathbb D_{i,X}^j
    \]
    is Kähler. Choose a Kähler form $\chi_i^j$ in this class and set $S_i^j=T_i^j+\chi_i^j$. Then $S_i^j$ lies in the fixed class $\mathbb D_{i,X}+\epsilon\{\omega\}$ and has the same divisorial Lelong numbers as $T_i^j$. Hence, by \cref{thm:usc_mix_vol} and \itemref{itm:prop:vollinearlimit:2},
    \[
        \varlimsup_{j\in J}\vol\Bigl(T_1^j,\ldots,T_n^j\Bigr)\leq \varlimsup_{j\in J}\vol\Bigl(S_1^j,\ldots,S_n^j\Bigr)\leq\vol\left(T_1+\epsilon\omega,\ldots,T_n+\epsilon\omega\right).
    \]
    Therefore,
    \[
        \varlimsup_{j\in J}\Bigl(\mathbb D_1^j,\ldots,\mathbb D_n^j\Bigr)\leq \Bigl(\mathbb D_1+\epsilon\mathbb D(\omega),\ldots,\mathbb D_n+\epsilon\mathbb D(\omega)\Bigr).
    \]
    Letting $\epsilon\to 0+$ gives \eqref{eq:limsupint_prod}.

    \textbf{Step~2}. Next we handle the general case.

    Take a Kähler form $\omega$ on $X$. Then by \cref{lma:perturbDF}, for any $\epsilon\in (0,1]$, we have
    \[
    \begin{aligned}
          \varlimsup_{j\in J}\left(\mathbb{D}_1^j,\ldots,\mathbb{D}_n^j \right)
        \leq & \varlimsup_{j\in J}\left(\mathbb{D}_1^j+\epsilon \mathbb{D}(\omega),\ldots,\mathbb{D}_n^j+\epsilon \mathbb{D}(\omega) \right)\\
        \leq & \left( \mathbb{D}_1+\epsilon \mathbb{D}(\omega),\ldots,\mathbb{D}_n+\epsilon \mathbb{D}(\omega) \right)\\
        \leq & \left( \mathbb{D}_1,\ldots,\mathbb{D}_n \right)+C\epsilon.
    \end{aligned}
    \]
    Since $\epsilon$ is arbitrary, our assertion follows.
\end{proof}

\begin{proposition}\label{prop:DFdecnet}
    The product in \cref{def:DF_ext} is continuous along decreasing nets in each variable. In other words, if $(\mathbb{D}_i^j)_{j\in J}$ ($i=1,\ldots,n$) are decreasing nets of nef b-divisors over $X$ with limits $\mathbb{D}_i$, then
    \[
    \lim_{j\in J} \left(\mathbb{D}_1^j,\ldots,\mathbb{D}_n^j\right)=\left(\mathbb{D}_1,\ldots,\mathbb{D}_n\right).
    \]
\end{proposition}
\begin{proof}
    This is a straightforward consequence of \itemref{itm:lma:DFlinear:4} and \cref{prop:DFusc}.
\end{proof}

\begin{remark}\label{rmk:bdiv:L749}
The algebraic nef b-divisors are introduced by Dang--Favre \cite{DF20}, as their intersection theory. It is a straightforward application of the results proved in this section that our transcendental intersection theory coincides with theirs in the algebraic setting.
\end{remark}
